\section{The Geometric Mass Gap}
\label{sec:geometric_gap}

This module provides an independent, geometric proof of the gap $\Delta > 0$ that relies on the curvature of the gauge orbit space $\mathcal{M} = \mathcal{A}/\mathcal{G}$, rather than functional inequalities.

\subsection{The O'Neill Curvature Formula}
The configuration space $\mathcal{M}$ is a Riemannian submersion of the flat connection space $\mathcal{A}$. By O'Neill's formula, the Ricci curvature of the quotient is:
\begin{equation}
\text{Ric}_{\mathcal{M}}(X, Y) = \text{Ric}_{\mathcal{A}}(X^h, Y^h) + 2 \langle A_X Y^h, A_Y X^h \rangle + \langle A_X X^h, A_Y Y^h \rangle
\end{equation}
where $A$ is the O'Neill tensor related to the vertical projection of the Lie bracket. For non-Abelian gauge groups ($[T^a, T^b] \neq 0$), the commutator is non-zero, yielding positive curvature:
\begin{equation}
\text{Ric}_{\mathcal{M}} \ge K > 0
\end{equation}

\subsection{The Lichnerowicz Spectral Gap}
By the Lichnerowicz theorem extended to infinite dimensions, a lower bound on Ricci curvature implies a spectral gap for the Laplacian (Hamiltonian):
\begin{equation}
\Delta \ge \frac{d}{d-1} K
\end{equation}
This proves $\Delta > 0$ purely from the non-commutative geometry of the gauge group, providing a geometric "safety net" for the gap.
